% ===== 7. Limitations & future work (DRAFTABLE) =============================
\section{Limitations and future work}
\label{sec:limitations}

\paragraph{Synthetic targets.}
All experiments use analytic shapes rendered into multi-view scenes, chosen so
the ground-truth surface and its Betti numbers are exact and so geometry can be
held fixed while topology varies---the control that makes the blindness result
and the topological-vs-non-topological comparison clean. The cost is external
validity: we do not yet demonstrate the method on captured or noisy real scenes,
where the target persistence diagram itself must be estimated.
% TODO(human): if you run any real/captured scene, add it; otherwise keep this honest.

\paragraph{Non-determinism of the optimizer.}
The CUDA rasterizer's gradient atomics make vertex positions non-reproducible
run-to-run; resampling structure ($F$) is deterministic but vertex positions
($V$) are not, so single-run differences can include a few drifting triangles. We
mitigate with paired seeds and seed-averaging and a soup-sampled topology metric,
but absolute numbers carry this run-to-run variance.
% src(determinism caveat): D:\Project\CG-Soup-Topology\PHASE2_STATUS.md

\paragraph{What the spread scale varies (limited dynamic range).}
For the compact targets used here, topological death scales are of order the shape
size, so even the concentrated field ($s{=}1$) already overlaps most of the surface.
The scale $s$ therefore controls the \emph{peakedness} of the importance field---and
thus the differential protection applied by the keep-map---rather than its literal
spatial coverage; our claims about ``spreading'' should be read in that sense.
Characterizing the effect on targets whose features are small relative to the shape,
where coverage and peakedness decouple, is left to future work.
% src(limited dynamic range / peakedness): PHASE2B_CROSSOVER.md (Caveats)

\paragraph{Width vs.\ topological placement.}
The void gain is primarily a kernel-width effect: a width-matched non-topological
control recovers most of it, and the residual attributable to topological placement
is small and shape-dependent---clear for the cylinder, modest for the sphere, absent
for the cube (Section~\ref{sec:results-h2}). Establishing that residual robustly, and
whether it grows for targets with finer or off-centre voids, would need more shapes
and seeds.
% src(width-matched control B5): D:\Project\CG-Soup-for-Digital-Dentistry\output\synth\h2_unified\report\results.json

\paragraph{Precomputed, non-differentiable prior.}
The importance field is distilled once from the target and is \emph{not}
differentiated through; the method biases sampling, it does not flow a topological
gradient. Making persistent homology differentiable so the topological signal can
backpropagate into vertex positions---rather than only steer resampling---is the
natural next phase, building on existing differentiable-persistence
machinery~\cite{gabrielsson2020topologylayer,carriere2021optimizing,poulenard2018shapematching}.
% src(scope: precomputed field, no differentiable PH): methods/topo_importance.py header

\paragraph{Resolving multi-loop / genus-2 targets.}
Our genus-2 \texttt{double\_torus} stress test was inconclusive: at every feasible
budget the soup recovered only 2 of its 4 loops and the decisive-dimension
bottleneck did not discriminate between conditions
(Section~\ref{sec:results-crossover}). Strengthening the spread-vs-concentrate
finding on a richer multi-loop target needs a higher-resolution or longer-schedule
re-run, which we leave to future work.
% src(double_torus inconclusive): D:\Project\CG-Soup-for-Digital-Dentistry\output\synth\crossover\crossover_report\results.json
